\documentclass[12pt]{article}
\usepackage[utf8]{inputenc}
\usepackage[T1]{fontenc}
\usepackage[margin=1in]{geometry}
\usepackage{setspace}
\usepackage{xcolor}
\usepackage[colorlinks=true, urlcolor=blue!60!black, linkcolor=blue!60!black]{hyperref}
\setlength{\parindent}{0em}
\setlength{\parskip}{0.6em}
\renewcommand{\baselinestretch}{1.0}
\begin{document}
\begin{center}
{\LARGE \textit{United States v. Hemani}, No. 24-1234 (2026)\par}
\vspace{0.5em}
{\large Case Brief --- Two-Page Edition\par}
\end{center}

\section{Memory Jogger}
A unanimous Court held that the Second Amendment bars prosecuting a regular marijuana user for keeping a gun at home under 18 U.S.C. \S~922(g)(3) on his drug use alone (\href{https://www.supremecourt.gov/opinions/25pdf/24-1234_g2bh.pdf}{Slip Op.}).

\section{Detailed Case Facts}
Agents searching Ali Hemani's family home in 2022 found a pistol and marijuana; Hemani admitted using marijuana ``about every other day.'' He was indicted only under \S~922(g)(3), which disarms any ``unlawful user of or addicted to any controlled substance.'' The government did not allege addiction, intoxication while armed, or dangerousness.

\section{Procedural History}
The district court dismissed the indictment; the Fifth Circuit affirmed under its \textit{Connelly} precedent. The Supreme Court granted review and decided the case on June 18, 2026.

\section{Judicial Votes}
Affirmed, 9--0. Gorsuch for seven (Roberts, Thomas, Sotomayor, Kavanaugh, Barrett, Jackson); Thomas and Jackson (with Sotomayor) concurring; Alito, with Kagan, concurring in the judgment.

\section{Holding}
Under \textit{Bruen} and \textit{Rahimi}, the government may not disarm someone under the ``unlawful user'' provision solely for regular drug use, without pre-deprivation process or proof of incapacitation or dangerousness. Laws aimed at addicts, the presently intoxicated, and individually dangerous persons were expressly reserved.

\section{Analysis of Opinions}
\textbf{Majority:} The government's ``habitual drunkard'' analogues failed on every metric. They targeted the persistently incapacitated rather than regular drinkers, served different purposes (labor norms, family protection, moral policing), and required judicial process before any deprivation. \textbf{Thomas:} \S~922(g) may exceed the commerce power. \textbf{Jackson:} the Court should retire \textit{Bruen} for means-end scrutiny. \textbf{Alito:} affirm narrowly, because the record showed nothing about how Hemani's use impaired him.

\section{Examples: Future Applications}
\textbf{Same-Side:} A medical-marijuana user with a locked shotgun at home and no violent history cannot be prosecuted on her regular use alone. \textbf{Opposite-Side:} A user who is actively high, armed, and threatening bystanders falls within \textit{Hemani}'s reservation for intoxication and proven dangerousness. \textbf{Fence-Sitter:} A frequent methamphetamine user arrested with a pistol, with no evidence his use made him violent or incapacitated, presents an open question.

\section{Critique}
Originalists praised the refusal to accept the government's bare ``say-so'' on dangerousness, while others faulted the Court for leaving little guidance on ``unlawful user.'' Following Justice Jackson, critics argue the decision reveals \textit{Bruen}'s incoherence and invites case-by-case fights over drug potency.

\section{Key Quotations}
\begin{quote}
``All of which leaves it awkwardly positioned to suggest that the millions of Americans who now regularly use marijuana are categorically and unusually dangerous.''
\end{quote}
\noindent --- Justice Neil M. Gorsuch, Majority Opinion, slip op. at 17.

\end{document}
